\documentclass[11pt]{article}
\usepackage[margin=1in]{geometry}
\usepackage{graphicx}
\usepackage{subcaption}
\usepackage{booktabs}
\usepackage{float}
\usepackage{hyperref}

\graphicspath{{figures/}}

% \onefig{file}{caption}
\newcommand{\onefig}[2]{%
  \begin{figure}[H]
    \centering
    \includegraphics[width=0.6\textwidth]{#1}
    \caption{#2}
\end{figure}}

% \twofig{left file}{left caption}{right file}{right caption}{overall caption}
\newcommand{\twofig}[5]{%
  \begin{figure}[H]
    \centering
    \begin{subfigure}{0.48\textwidth}
      \includegraphics[width=\linewidth]{#1}
      \caption{#2}
    \end{subfigure}\hfill
    \begin{subfigure}{0.48\textwidth}
      \includegraphics[width=\linewidth]{#3}
      \caption{#4}
    \end{subfigure}
    \caption{#5}
\end{figure}}

% \tablerow{file}: pulls in one table row (plain \input breaks inside a tabular)
\makeatletter
\newcommand{\tablerow}[1]{\@@input tables/#1.tex }
\makeatother

% \tableinput{file}{caption}
\newcommand{\tableinput}[2]{%
  \begin{table}[H]
    \centering
    \caption{#2}
    \input{tables/#1}
\end{table}}

\title{CS 4372 Assignment 2: Taiwanese Bankruptcy Prediction}
\author{}
\date{\today}

\begin{document}
\maketitle

% =====================================================================
\section{Data}

The data set chosen for this assignment is the Taiwanese Bankruptcy
Prediction data set from the UC Irvine Machine Learning Repository.
The data was collected from 1999 to 2009 and includes a little under
7000 companies. The target variable is if the company defaulted, and
there are 95 continuous features which cover various aspects of a
companies balance sheet, such as revenue per employee, realized sales
gross margin, and various ratios. A tree is a good fit for a dataset
such as this since many of the feature variables would be highly
correlated with one another, making regression difficult. For
example, the cash assets to total assets ratio would be colinear with
the cash to current liability ratio.

One challenge in this data set is that, by far, the vast majority of
data points have a feature variable equal to 0, that is, they did not default.

By the nature of this data set, we are mostly interested in being
able to predict if a company would fail based on their current
balance sheet, as such, we are mostly interested in the recall and F1
statistic of a model. Through the report, the predictive power of the
model will be emphesized.

% =====================================================================
\section{Methodology}

All four models were fit with the same procedure, implemented in
scikit-learn (and the \textit{xgboost} package for XGBoost).

\paragraph{Train/test split.} The data was split once into 80\%
training and 20\% test partitions, stratified on the target so that
both partitions keep the same proportion of bankrupt companies. The
same split (fixed random seed) is used for every model, so all test-set
results are computed on identical data. The test partition is not used
for any fitting or tuning.

\paragraph{Baseline model.} Each model was first fit on the training
partition with its library default parameters, and its
precision--recall curve on the test partition was recorded as a
baseline.

\paragraph{Grid search.} Hyperparameters were selected with an
exhaustive grid search using 5-fold stratified cross-validation on the
training partition, scored by F1 on the bankrupt (positive) class.
Training-fold scores were also recorded. The parameter combination
with the highest mean validation F1 was then refit on the full
training partition; this refit model is the final estimator for that
model family.

\paragraph{Validation curves.} To isolate the effect of a single
complexity parameter (\textit{max\_depth} for the tree-based models,
\textit{n\_estimators} for the boosting models), that parameter was
varied over a range while all others were left at their defaults. At
each value the mean and standard deviation of the training and
validation F1 over 5-fold cross-validation are reported. A second plot
uses the grid-search results instead, showing for each value of the
same parameter the best mean F1 over all combinations of the remaining
grid parameters.

\paragraph{Threshold tuning.} The final estimator predicts the
bankrupt class when its predicted probability exceeds 0.5. Because the
classes are heavily imbalanced, a second decision threshold was
chosen: the final estimator was wrapped in scikit-learn's
\textit{TunedThresholdClassifierCV}, which evaluates candidate
thresholds by 5-fold cross-validation on the training partition,
selects the threshold maximizing F1, and refits the model on the full
training partition.

\paragraph{Evaluation.} The final estimator was evaluated on the test
partition at both the 0.5 threshold and the tuned threshold, reporting
accuracy, balanced accuracy, precision, recall, F1, and the Matthews
correlation coefficient (MCC), along with the confusion matrix at each
threshold. ROC AUC and average precision (the area under the
precision--recall curve) are computed from the predicted probabilities
and do not depend on the threshold.

\section{Decision Tree}

The first model fit was a naive descision tree fit with the default
parameters from scikit learn. As expected, since the tree fits until
all samples in leaf nodes are pure, the precision/recall curve has
one point somewhere around the middle of the graph (at a 0.5
threshold) where it is not either all precision or all recall. A grid
search with 5 fold cross validation over the \textit{ccp\_alpha} and
\textit{max\_depth} using F1 as the scoring criterion was conducted.

The candidate \textit{ccp\_alpha} values were taken from the
cost-complexity pruning path of the unpruned tree on the training
partition, i.e.\ the values of $\alpha$ at which pruning removes
another subtree. The largest value, which prunes the tree down to a
single root node, was dropped, and every tenth remaining value was
kept to limit the size of the grid. The candidate \textit{max\_depth}
values were every integer from 1 up to (but not including) the depth
of the unpruned tree, plus an unrestricted depth. All trees use the
Gini impurity criterion and consider every feature at each split. The
validation curve varies \textit{max\_depth} from 1 to the depth of the
unpruned tree with no pruning ($\alpha = 0$).

\subsection{Hyperparameter Search}
\onefig{descision_tree_unpruned_pr_curve}{Precision--recall curve of
the unpruned decision tree (default parameters).}
\tableinput{descision_tree_best_params}{Decision tree: best
hyperparameters from the grid search.}
\twofig{descision_tree_f1_vs_max_depth}{Validation curve}%
{descision_tree_grid_f1_vs_max_depth}{Best grid-search F1 per depth
(over \texttt{ccp\_alpha})}%
{Decision tree: F1 vs.\ \texttt{max\_depth}.}

When regularization is introduced, the best CV F1 per max depth came out to
be 0.366 with a max depth of 12. Without any alpha, the best max
depth came out to be 8. The difference in explainability between 8
and 12 is relatively minimal, so we proceed with the best descision
tree that the grid search came out to.

\subsection{Test-Set Evaluation}
\tableinput{descision_tree_metrics}{Decision tree: test-set metrics
at the default and tuned thresholds.}
\twofig{descision_tree_confusion_matrix}{Threshold 0.5}%
{descision_tree_tuned_confusion_matrix}{Tuned threshold}%
{Decision tree: test-set confusion matrices.}
\twofig{descision_tree_pr_curve}{Precision--recall curve}%
{descision_tree_roc_curve}{ROC curve}%
{Decision tree: test-set precision--recall and ROC curves.}

As is the case with the rest of the models in this report, the ROC
AUC is high since the vast majority of the dataset did not default,
as such, the model will get the vast majority of its classifications
correct, but we only care about the precision and recall of the model.

% =====================================================================
\section{Random Forest}

A random forest fits many decision trees, each on a bootstrap sample
of the training partition, considering a random subset of
$\sqrt{95} \approx 9$ features at each split, and predicts the average
of the trees' class probabilities. The baseline used scikit-learn's
defaults: 100 trees grown to purity with the Gini criterion.

The grid search covered \textit{n\_estimators} $\in \{100, 300\}$,
\textit{max\_depth} $\in \{5, 10, 20, \textrm{unrestricted}\}$, and
\textit{class\_weight} $\in \{\textrm{none}, \textrm{balanced}\}$, for 16
combinations. The balanced setting weights each class inversely to
its frequency in the training data, so errors on the rare bankrupt
class cost proportionally more when growing the trees. The validation
curve varies \textit{max\_depth} over $2, 4, \dots, 30$ with 100 trees
and no class weighting.

\subsection{Hyperparameter Search}
\onefig{random_forest_default_pr_curve}{Precision--recall curve of
the random forest with default parameters.}
\tableinput{random_forest_best_params}{Random forest: best
hyperparameters from the grid search.}
\twofig{random_forest_f1_vs_max_depth}{Validation curve}%
{random_forest_grid_f1_vs_max_depth}{Best grid-search F1 per depth
(over \texttt{n\_estimators}, \texttt{class\_weight})}%
{Random forest: F1 vs.\ \texttt{max\_depth}.}

The random forest model benefit much more from the regularization
parameter class weight as well as from being provided with additional
estimators. A mean F1 of 0.425 is a definite improvement which holds
up in the test set.

\subsection{Test-Set Evaluation}
\tableinput{random_forest_metrics}{Random forest: test-set metrics at
the default and tuned thresholds.}
\twofig{random_forest_confusion_matrix}{Threshold 0.5}%
{random_forest_tuned_confusion_matrix}{Tuned threshold}%
{Random forest: test-set confusion matrices.}
\twofig{random_forest_pr_curve}{Precision--recall curve}%
{random_forest_roc_curve}{ROC curve}%
{Random forest: test-set precision--recall and ROC curves.}

Although worse on the precision and about the same in accuracy, the
random forest model outperformed the descision tree in recall. This
is beneficial since we are more tolerant of a false positive than we
are of a false negative (that is, the possible financial loss is
  larger from a bankruptcy we failed to predict than incorrectly
diagnosing a company as high risk).

% =====================================================================
\section{AdaBoost}

AdaBoost fits a sequence of decision stumps (depth-1 trees). After
each round, the sample weights of
misclassified training points are increased so the next stump focuses
on them, and each stump receives a vote weight based on its weighted
error. The prediction is the weighted vote of all stumps. The
baseline used the defaults of 50 stumps and a learning rate of 1.0;
the learning rate shrinks each stump's vote weight.

The grid search covered \textit{n\_estimators} $\in \{50, 100, 200,
400\}$ and \textit{learning\_rate} $\in \{0.1, 0.5, 1.0\}$, for 12
combinations. The validation curve varies \textit{n\_estimators} over
$\{25, 50, 100, 200, 300, 400, 600\}$ at a learning rate of 1.0.

For the boosting diagnostics, the training partition was split again
into 80\% fitting and 20\% validation subsets (stratified). A model
with the best grid-search parameters was fit on the fitting subset,
and the average precision on both subsets was computed after every
boosting round. The per-round stump weights and weighted errors, and
the feature importances (each stump's impurity-based importance,
averaged with the stump weights), are taken from the final estimator.

\subsection{Hyperparameter Search}
\onefig{adaboost_default_pr_curve}{Precision--recall curve of
AdaBoost with default parameters.}
\tableinput{adaboost_best_params}{AdaBoost: best hyperparameters from
the grid search.}
\twofig{adaboost_f1_vs_n_estimators}{Validation curve}%
{adaboost_grid_f1_vs_n_estimators}{Best grid-search F1 per
\texttt{n\_estimators} (over \texttt{learning\_rate})}%
{AdaBoost: F1 vs.\ \texttt{n\_estimators}.}
\onefig{adaboost_learning_rate_vs_n_estimators}{AdaBoost:
cross-validated F1 vs.\ \texttt{n\_estimators} for each learning rate.}

The number of estimators agreed between training with the learning
rate and without, which indicates that the learning rate is not
particularly meaningful with this data set. This is confirmed when we
plot the learning rate against number of estimators and the CV F1.

\subsection{Test-Set Evaluation}
\tableinput{adaboost_metrics}{AdaBoost: test-set metrics at the
default and tuned thresholds.}
\twofig{adaboost_confusion_matrix}{Threshold 0.5}%
{adaboost_tuned_confusion_matrix}{Tuned threshold}%
{AdaBoost: test-set confusion matrices.}
\twofig{adaboost_pr_curve}{Precision--recall curve}%
{adaboost_roc_curve}{ROC curve}%
{AdaBoost: test-set precision--recall and ROC curves.}
% Analysis

\subsection{Boosting Diagnostics}
\twofig{adaboost_staged_average_precision}{Average precision per
boosting round}%
{adaboost_estimator_weights_errors}{Per-round estimator weights and
weighted errors}%
{AdaBoost: behaviour across boosting rounds.}
\onefig{adaboost_feature_importance}{AdaBoost: top 15 features by importance.}

After tuning the threshold, AdaBoost does about as well as the random
forest model does. I suspect that if we biased the adaptive boosting
weights towards towards giving false negatives higher weight than
false positives, we would end up with a model which more completly
identifies bankrupting companies.

% =====================================================================
\section{XGBoost}

XGBoost fits gradient-boosted trees: each round adds a tree fit to the
gradient of the logistic loss of the current ensemble, scaled by the
learning rate, with regularization on leaf weights. The baseline used
the library defaults of 100 rounds, a learning rate of 0.3, and a
maximum tree depth of 6, with histogram-based split finding.

The grid search covered \textit{n\_estimators} $\in \{100, 200,
400\}$, \textit{max\_depth} $\in \{3, 6\}$, \textit{learning\_rate}
$\in \{0.1, 0.3\}$, and \textit{scale\_pos\_weight} $\in \{1, w\}$,
for 24 combinations, where $w$ is the ratio of non-bankrupt to
bankrupt companies in the training partition. Setting
\textit{scale\_pos\_weight} to $w$ multiplies the loss on bankrupt
companies by $w$, so both classes contribute equally in total. The
validation curve varies \textit{n\_estimators} over $\{25, 50, 100,
200, 300, 400, 600\}$ with all other parameters at their defaults.

For the boosting diagnostics, the training partition was split again
into 80\% fitting and 20\% validation subsets (stratified). A model
with the best grid-search parameters, but allowed up to 1000 rounds,
was fit on the fitting subset while the area under the
precision--recall curve (AUCPR) was tracked on both subsets after
every round. Training stopped once the validation AUCPR had not
improved for 50 consecutive rounds, and the best round is reported.
Feature importance is measured by gain: the average reduction in loss
from the splits that use each feature in the final estimator.

\subsection{Hyperparameter Search}
\onefig{xgboost_default_pr_curve}{Precision--recall curve of XGBoost
with default parameters.}
\tableinput{xgboost_best_params}{XGBoost: best hyperparameters from
the grid search.}
\twofig{xgboost_f1_vs_n_estimators}{Validation curve}%
{xgboost_grid_f1_vs_n_estimators}{Best grid-search F1 per
\texttt{n\_estimators} (over the other parameters)}%
{XGBoost: F1 vs.\ \texttt{n\_estimators}.}
\onefig{xgboost_learning_rate_vs_n_estimators}{XGBoost:
cross-validated F1 vs.\ \texttt{n\_estimators} for each learning rate.}
The grid search preferred the slower configuration: a learning rate
of 0.1, shallow trees of depth 3, the largest number of rounds in the
grid (400), and \textit{scale\_pos\_weight} set to $w \approx 30$.
XGBoost was the only boosting model given a way to weight the
bankrupt class directly, and the grid search chose to use it.

\subsection{Test-Set Evaluation}
\tableinput{xgboost_metrics}{XGBoost: test-set metrics at the default
and tuned thresholds.}
\twofig{xgboost_confusion_matrix}{Threshold 0.5}%
{xgboost_tuned_confusion_matrix}{Tuned threshold}%
{XGBoost: test-set confusion matrices.}
\twofig{xgboost_pr_curve}{Precision--recall curve}%
{xgboost_roc_curve}{ROC curve}%
{XGBoost: test-set precision--recall and ROC curves.}
At the default threshold XGBoost had the highest F1 (0.484) and MCC
(0.466) of the four models, mostly because of its precision (0.468).
Tuning lowered the threshold to 0.414, which traded a small amount of
precision for recall (0.500 to 0.568) and raised F1 to 0.500. Even so,
XGBoost still identifies fewer of the 44 bankrupt companies in the
test set (25) than the random forest does (29).

\subsection{Boosting Diagnostics}
\twofig{xgboost_staged_aucpr}{AUCPR per boosting round (early stopping)}%
{xgboost_feature_importance}{Top 15 features by gain}%
{XGBoost: behaviour across boosting rounds and feature importance.}
% Analysis

% =====================================================================
\section{Model Comparison}

\begin{table}[H]
  \centering
  \caption{Test-set results for all models. ``Tuned'' uses the
  F1-maximizing threshold chosen by cross-validation on the training set.}
  \begin{tabular}{lcccccc}
    \toprule
    & \multicolumn{2}{c}{F1} & \multicolumn{2}{c}{MCC} & & \\
    \cmidrule(lr){2-3} \cmidrule(lr){4-5}
    Model & 0.5 & Tuned & 0.5 & Tuned & ROC AUC & Avg.\ precision \\
    \midrule
    \tablerow{descision_tree_summary_row}
    \tablerow{random_forest_summary_row}
    \tablerow{adaboost_summary_row}
    \tablerow{xgboost_summary_row}
    \bottomrule
  \end{tabular}
\end{table}
All three ensembles clearly outperformed the single decision tree,
which had the lowest ROC AUC (0.827) and average precision (0.375).
Among the ensembles the differences are small: with only 44 bankrupt
companies in the test set, a difference of 0.02--0.04 in F1 amounts to
two or three companies. XGBoost has the best F1 and MCC, but the
random forest has the highest recall (0.659), balanced accuracy
(0.809), ROC AUC (0.949), and average precision (0.464), and it reaches
these at the default threshold without any tuning, since its
\textit{class\_weight} setting already accounts for the imbalance.
Because missing a bankruptcy is more costly than flagging a healthy
company, recall matters most for this task, so the random forest is
the best of the four models here.

\end{document}
